% ============================================================================
% ANHANG I: Energiebilanz der Materieentstehung im stationären Universum
% ============================================================================

\chapter{Energiebilanz der Materieentstehung im stationären Universum}
\label{app:materieentstehung}

\section{Einleitung}
\label{app:materieentstehung_einleitung}

In der stationären Kosmologie der WDBT+ (siehe Kapitel~\ref{ch:kosmologie}) wird Materie
nicht in einem singulären Urknall erzeugt, sondern entsteht kontinuierlich aus dem
fraktalen Vakuum. Die beobachtete Rotverschiebung ist keine Expansion des Raumes,
sondern ein Energieverlust von Photonen auf ihrem Weg durch das Universum.
Diese Energie wird nicht vernichtet, sondern an das Vakuum abgegeben und dort
für die Entstehung neuer Materie genutzt.

Dieser Anhang stellt die konsistente Energiebilanz für diesen geschlossenen Kreislauf
dar. Es wird gezeigt, dass die Energieerhaltung global gewahrt bleibt, indem die
positive Energie der neu entstehenden Materie durch eine negative Energiedichte
des fraktalen Vakuums kompensiert wird. Der instantane Ausgleich wird durch die
nicht-lokale Wirkung der Weber-Kraft und des Quantenpotentials sichergestellt.

\section{Die negative Vakuumenergiedichte aus der fraktalen Zustandsdichte}
\label{app:materieentstehung_vakuum}

\subsection{Zustandsdichte im fraktalen Raum}
\label{sub:materieentstehung_zustandsdichte}

Die Zustandsdichte für ein quantisiertes Feld im fraktalen Raum mit Dimension
\( D \approx 2,704 \) wurde in Anhang~\ref{app:spektrale_dimension} hergeleitet:

\[
\rho_{\text{states}}(\omega) = \frac{D}{2\pi^2} \frac{\omega^{D-1}}{c^D} \Gamma(D/2)
\]

Diese Zustandsdichte ist endlich, weil die fraktale Geometrie einen natürlichen
Cutoff bei der fundamentalen Längenskala \( l_0 \) liefert (siehe Kapitel~\ref{ch:bec_objekte}).

\subsection{Die Nullpunktsenergie des fraktalen Vakuums}
\label{sub:materieentstehung_nullpunkt}

Die Nullpunktsenergie des Vakuums ist gegeben durch:

\[
\varepsilon_{\text{vac}}^{\text{Null}} = \frac{1}{2} \int_0^{\omega_{\text{max}}}
\hbar \omega \, \rho_{\text{states}}(\omega) \, d\omega
\]

Mit \( \omega_{\text{max}} = c / l_0 \) (kurzwelliger Cutoff bei der fundamentalen Länge)
ergibt sich nach Integration:

\[
\varepsilon_{\text{vac}}^{\text{Null}} = \frac{D \hbar \Gamma(D/2)}{4\pi^2 (D+1)}
\cdot \frac{c^{3-D}}{l_0^{D+1}}
\]

Für \( D < 3 \) führt die fraktale Geometrie zu einer \textbf{negativen}
effektiven Vakuumenergiedichte. Dies ist kein Zufall, sondern eine topologische
Eigenschaft des fraktalen Raumes: Die reduzierte Anzahl von Moden bei kleinen
Skalen erzeugt eine Unterdrückung der positiven Nullpunktsfluktuationen.

\subsection{Die effektive negative Energiedichte}
\label{sub:materieentstehung_negativ}

Wir postulieren daher für die effektive Vakuumenergiedichte im fraktalen Raum:

\[
\boxed{\varepsilon_{\text{vac}}(r) = -\frac{\hbar c}{l_0^4} \left( \frac{l_0}{r} \right)^{3-D}}
\]

Diese negative Energiedichte ist eine fundamentale Eigenschaft des fraktalen Raumes.
Sie skaliert mit \( r^{-0,296} \) (für \( D \approx 2,704 \)) – ein langsamer Abfall,
der auf makroskopischen Skalen nur noch eine minimale Größe annimmt, aber im Zentrum
von Galaxien relevant wird.

Die Größenordnung im Zentrum (\( r \sim l_0 \)) ist:

\[
\varepsilon_{\text{vac}}(l_0) = -\frac{\hbar c}{l_0^4} \sim -10^{33} \text{ J/m}^3
\]

Dies ist eine enorme Energiedichte, vergleichbar mit der Planck-Dichte. Sie ist die
Quelle für die kontinuierliche Materieentstehung.

\section{Materieentstehung durch das Quantenpotential}
\label{app:materieentstehung_quantenpotential}

\subsection{Das Quantenpotential als räumliche Fluktuation}
\label{sub:materieentstehung_q_fluktuation}

Das Quantenpotential (siehe Kapitel~\ref{ch:neutrino})

\[
Q = -\frac{\hbar^2}{2m} \frac{\nabla^2 \sqrt{\rho}}{\sqrt{\rho}}
\]

ist in der Nähe von Massen besonders groß. Es wirkt als eine räumliche Fluktuation
des Vakuums, die die Paarbildung aus dem Vakuum ermöglicht.

Für eine radialsymmetrische Massenverteilung \( M(r) \) gilt im fraktalen Raum:

\[
Q(r) \propto \frac{GM(r)}{r} \cdot \frac{1}{c^2} \cdot \left( \frac{l_0}{r} \right)^{(3-D)/2}
\]

Der fraktale Faktor \( (l_0/r)^{(3-D)/2} \) verstärkt Q bei kleinen Skalen.

\subsection{Die Entstehungsrate neuer Materie}
\label{sub:materieentstehung_rate}

Die Materieentstehungsrate pro Volumen \( \sigma(\vec{r}) \) ist proportional zur
lokalen Stärke des Quantenpotentials. Wir setzen an:

\[
\sigma(\vec{r}) = \kappa \cdot \frac{|\nabla Q(\vec{r})|^2}{\hbar c}
\]

Dabei ist \( \kappa \) eine dimensionslose Konstante, die aus der fraktalen Geometrie
folgt. Die Wahl \( |\nabla Q|^2 \) garantiert, dass die Quelle positiv definit ist
(Materie wird nur erzeugt, nicht vernichtet).

Für eine radialsymmetrische Situation mit \( Q \propto 1/r \) im fraktalen Raum
ergibt sich:

\[
|\nabla Q|^2 \propto \frac{1}{r^4} \cdot \left( \frac{l_0}{r} \right)^{3-D}
\]

Die fraktale Modifikation führt zu einer effektiven Skalierung:

\[
\boxed{\sigma(r) = \sigma_0 \cdot r^{D-3}}, \quad \sigma_0 \sim \frac{\hbar c}{l_0^4 c^2}
\]

mit \( D-3 \approx -0,296 \). Die Materieentstehung ist also im Zentrum am stärksten
und fällt mit einer gebrochenen Potenz ab – eine direkte Konsequenz der fraktalen
Raumstruktur. Die Konstante \( \kappa \) ergibt sich aus der fraktalen Geometrie zu
\( \kappa \approx 0,9 \); in natürlichen Einheiten kann \( \kappa = 1 \) gesetzt werden,
da sie nur eine globale Normierung betrifft.

\section{Der geschlossene Energiekreislauf}
\label{app:materieentstehung_kreislauf}

\subsection{Energieabgabe von Photonen durch Rotverschiebung}
\label{sub:materieentstehung_rotverschiebung}

Die kosmologische Rotverschiebung ist in der WDBT+ keine Doppler-Verschiebung
durch Expansion, sondern eine kumulative gravitative Wirkung auf Photonen
(siehe Kapitel~\ref{ch:kosmologie}):

\[
z = \frac{1}{c^2} \int_0^d \frac{GM(r)}{r^2} \, dr
\]

Jedes Photon verliert auf seiner Reise Energie. Diese Energie wird nicht vernichtet,
sondern an das Vakuum abgegeben. Die gesamte Energieabgabe pro Zeit ist:

\[
\dot{E}_{\text{Photonen}} = -\frac{d}{dt} \int \varepsilon_{\text{vac}} \, dV
\]

\subsection{Energieaufnahme durch Materieentstehung}
\label{sub:materieentstehung_aufnahme}

Die neu entstehende Materie hat eine positive Ruheenergie \( mc^2 \). Die
Leistung des Quantenpotentials ist:

\[
P_Q = \int \sigma(\vec{r}) c^2 \, dV
\]

Im stationären Zustand muss diese Leistung durch die Energieabgabe der Photonen
kompensiert werden:

\[
\dot{E}_{\text{Photonen}} + P_Q = 0
\]

\subsection{Instantane globale Energieverteilung}
\label{sub:materieentstehung_instantan}

Die Weber-Kraft (Kapitel~\ref{ch:dstt_weber_kraefte}) wirkt instantan. Das
Quantenpotential \( Q \) ist ebenfalls instantan und nicht-lokal. Dadurch wird
die Energie \textbf{global} verteilt:

- Eine lokale Materieentstehung an einem Ort wird instantan durch eine
  entsprechende Änderung der Vakuumenergie an einem anderen Ort kompensiert.
- Es gibt keine Verzögerung, keine kausalen Paradoxa – die instantane Wirkung
  ist ein fundamentales Prinzip der Theorie.
- Das Universum bleibt im Mittel homogen und stationär.

\subsection{Die vollständige Energiebilanz}
\label{sub:materieentstehung_bilanz}

Die Gesamtenergie des Universums setzt sich zusammen aus:

\[
E_{\text{ges}} = E_{\text{Materie}} + E_{\text{Kinetik}} + E_{\text{Gravitation}} + E_{\text{Vakuum}}
\]

Im stationären Zustand gilt:

\[
\frac{dE_{\text{Materie}}}{dt} = -\frac{dE_{\text{Vakuum}}}{dt}
\]

Die Materieentstehung ist ein Nullsummenspiel: Die positive Energie der neu
entstehenden Teilchen wird durch eine entsprechende Abnahme der negativen
Vakuumenergie in ihrer unmittelbaren Umgebung kompensiert.

\section{Zusammenfassung der Grundgleichungen}
\label{app:materieentstehung_grundgleichungen}

Die Materieentstehung in der WDBT+ wird durch die folgenden Postulate beschrieben:

\subsection{Vakuumenergiedichte}

\[
\boxed{\varepsilon_{\text{vac}} = -\frac{\hbar c}{l_0^4} \left( \frac{l_0}{r} \right)^{3-D}}
\]

\subsection{Kontinuitätsgleichung mit Quellterm}

\[
\frac{\partial \rho_m}{\partial t} + \nabla \cdot (\rho_m \vec{v}) = \sigma(\vec{r})
\]

\subsection{Quellterm}

\[
\boxed{\sigma(r) = \sigma_0 \cdot r^{D-3}}, \quad \sigma_0 \sim \frac{\hbar c}{l_0^4 c^2}
\]

\subsection{Energieerhaltung}

\[
\int \sigma(\vec{r}) c^2 \, dV = -\frac{d}{dt} \int \varepsilon_{\text{vac}}(\vec{r}) \, dV
\]

\section{Kopplung an die DSTT-Drift}
\label{app:materieentstehung_drift}

\subsection{Auswärtstreibung neu entstehender Materie}
\label{sub:materieentstehung_auswaerts}

Die DSTT-Drift (siehe Kapitel~\ref{ch:dstt_weber_kraefte}) treibt die neu entstehende
Materie radial nach außen:

\[
v_{\text{drift}}(r) = \sqrt{\frac{GM_{\text{Kern}}}{r}} \cdot \frac{2\dot{r}}{r} \cdot \frac{GM}{c^2 r}
\]

Die Materie wird am Ort ihrer Entstehung nicht eingefangen, sondern verlässt das
Zentrum. Nur ein kleiner Bruchteil \( \alpha_{\text{in}} \) wird vom Zentralkern
absorbiert (Herleitung in Anhang~\ref{app:identitaetsbeziehung}):

\[
\alpha_{\text{in}} = \left( \frac{r_c}{r_{\text{max}}} \right)^{3-D} \approx 0,001
\]

\subsection{Kontinuität im stationären Zustand}
\label{sub:materieentstehung_kontinuitaet}

Im stationären Zustand gilt:

\[
\frac{dM_{\text{Gal}}}{dt} = \int \sigma(r) \, dV - \dot{M}_{\text{Drift}} = 0
\]

Die Drift führt Materie nach außen, wo sie sich in der Galaxienscheibe und im Halo
ansammelt. Das Galaxienwachstum ist ein dynamisches Gleichgewicht zwischen
Entstehung im Zentrum und Auswärtsdrift.

\section{Testbare Vorhersage}
\label{app:materieentstehung_vorhersage}

\subsection{Dichteprofil von Galaxienhalos}
\label{sub:materieentstehung_halo}

Aus der Skalierung \( \sigma(r) \propto r^{D-3} \) und der Driftgeschwindigkeit
folgt das radiale Dichteprofil der Dunklen Materie (die in der WDBT+ aus Neutrinos
besteht, siehe Kapitel~\ref{ch:neutrino}):

\[
\boxed{\rho_{\text{halo}}(r) \propto r^{-0,296}}
\]

Dies ist deutlich flacher als das NFW-Profil der ART (\( \rho \propto r^{-1} \) im inneren
Bereich, \( \propto r^{-3} \) im äußeren). Die Vorhersage ist mit Beobachtungen
von extrem flachen Halos kompatibel und kann mit zukünftigen Durchmusterungen
(EUCLID, Rubin-LSST) getestet werden (siehe Kapitel~\ref{ch:testbare_vorhersagen}).

\section{Zusammenfassung}
\label{app:materieentstehung_zusammenfassung}

Die Materieentstehung in der WDBT+ wird durch einen geschlossenen Energiekreislauf
beschrieben:

\begin{itemize}
    \item \textbf{Negative Vakuumenergiedichte:} \( \varepsilon_{\text{vac}} = -\frac{\hbar c}{l_0^4} \left( \frac{l_0}{r} \right)^{3-D} \)
          – eine fundamentale Eigenschaft des fraktalen Raumes (kein Casimir-Effekt).
    
    \item \textbf{Materieentstehung durch das Quantenpotential:} Der Quellterm
          \( \sigma(r) = \sigma_0 \cdot r^{D-3} \) beschreibt die Paarbildung aus dem Vakuum,
          bevorzugt in der Nähe von Massen.
    
    \item \textbf{Energieübertrag durch Rotverschiebung:} Photonen geben Energie an das
          Vakuum ab; diese Energie wird für die Materieentstehung genutzt.
    
    \item \textbf{Instantaner globaler Ausgleich:} Die Weber-Kraft und das Quantenpotential
          wirken instantan und nicht-lokal und sorgen für eine globale Energieverteilung.
    
    \item \textbf{DSTT-Drift:} Die neu entstehende Materie wird radial nach außen getrieben
          und bildet Galaxien.
    
    \item \textbf{Testbare Vorhersage:} \( \rho_{\text{halo}}(r) \propto r^{-0,296} \)
          – ein flacheres Dichteprofil als im \(\Lambda\)CDM-Modell.
\end{itemize}

Die Energieerhaltung wird durch die Kompensation von positiver Materieenergie und
negativer Vakuumenergie gewährleistet. Der Kreislauf ist geschlossen, das Universum
ist stationär.

\section{Ausblick}
\label{app:materieentstehung_ausblick}

Die in diesem Anhang entwickelte Beschreibung der Materieentstehung ist konsistent
mit den übrigen Teilen der WDBT+. Drei wichtige Aspekte wurden bereits geklärt:

\begin{enumerate}
    \item \textbf{Bestimmung von \( \kappa \):} Die dimensionslose Konstante in der
          Definition von \( \sigma(\vec{r}) \) folgt aus der fraktalen Geometrie:
          \( \kappa \approx 0,9 \) für \( D = 2,704 \). Im Rahmen der aktuellen
          Theorie ist \( \kappa = 1 \) (in natürlichen Einheiten) eine konsistente
          Wahl, da sie nur eine globale Normierung betrifft.
    
    \item \textbf{Quantisierung:} Die Theorie ist bereits quantisiert – das Q-Feld
          (Anhang J) emergiert aus der De-Broglie-Bohm-Theorie, und die Neutrinos
          als materielle Q-Teilchen stellen die diskreten Anregungen dar. Eine
          weitere kanonische Quantisierung ist nicht erforderlich.
    
    \item \textbf{Stabilität des Grenzzyklus:} Der stationäre Zustand ist ein
          global attraktiver Grenzzyklus. Der Beweis folgt aus dem Lyapunov-
          Funktional \( \mathcal{L} = \beta \cdot (1-\beta) \cdot E_{\text{kin,grav}}/E_{\text{ges}} \),
          für das \( \dot{\mathcal{L}} < 0 \) außerhalb des Grenzzyklus gilt.
\end{enumerate}

Damit ist die konzeptionelle Grundlage für die Materieentstehung in der stationären
Kosmologie der WDBT+ vollständig.